\documentclass{article}
\usepackage{amsmath}
\usepackage{graphicx}
\usepackage{geometry}
\geometry{a4paper, margin=1in}
\title{SelAction Program: Technical Description of \texttt{seldiscrete.f90}}
\author{Generated by Gemini}
\date{\today}

\begin{document}

\maketitle

\begin{abstract}
This document provides a detailed technical description of the \texttt{seldiscrete.f90} module from the SelAction program. This module implements one, two, and three-stage selection for populations with discrete generations.
\end{abstract}

\tableofcontents

\section{One-Stage Selection (\texttt{sel1s})}

The \texttt{sel1s} subroutine implements a single stage of selection.

\subsection{Selection Index}
The core of the calculation is the selection index, $I = \mathbf{b'x}$, where the vector of index weights, $\mathbf{b}$, is calculated as:
\begin{equation}
\mathbf{b} = \mathbf{P}^{-1}\mathbf{Gv}
\end{equation}
The response to selection for each trait is then calculated as:
\begin{equation}
\Delta G_i = \frac{\mathbf{b'G}_i}{\sigma_I} i
\end{equation}
where $\mathbf{G}_i$ is the $i$-th row of the $\mathbf{G}$ matrix, $\sigma_I$ is the standard deviation of the index, and $i$ is the selection intensity.

\subsection{Bulmer Equilibrium}
The program iteratively recalculates the genetic and phenotypic covariance matrices to account for the reduction in variance due to selection (the Bulmer effect). The additive genetic covariance matrix is updated as follows:
\begin{equation}
\mathbf{A}^* = \mathbf{A} - k \frac{(\mathbf{Ab})(\mathbf{Ab})'}{\mathbf{b'Pb}}
\end{equation}
where $\mathbf{A}$ is the additive genetic covariance matrix, $\mathbf{P}$ is the phenotypic covariance matrix, $\mathbf{b}$ is the vector of index weights, and $k$ is the variance reduction coefficient. This process is repeated for 25 iterations.

\section{Two-Stage Selection (\texttt{sel2s})}
The \texttt{sel2s} subroutine implements two-stage selection. The total genetic gain is the sum of the gains from each stage. The key is that the parameters for the second stage are conditional on the selection in the first stage.

\subsection{First Stage}
The first stage of selection is identical to the one-stage case, resulting in a genetic gain of $\Delta G_1$.

\subsection{Second Stage}
The second stage of selection is performed on the individuals selected in the first stage. The genetic and phenotypic covariance matrices are updated to account for the first stage of selection. The genetic gain from the second stage, $\Delta G_2$, is calculated using the updated parameters. The total genetic gain is $\Delta G_{total} = \Delta G_1 + \Delta G_2$.

\section{Three-Stage Selection (\texttt{sel3s})}
The \texttt{sel3s} subroutine extends the two-stage selection to three stages. The logic is analogous, with the parameters for each subsequent stage being updated based on the selection in the previous stages. The total genetic gain is:
\begin{equation}
\Delta G_{total} = \Delta G_1 + \Delta G_2 + \Delta G_3
\end{equation}

\end{document}
